% !PRINTED 7
\weberpage{7}
%<<CONT
ein Element $a$ zu $A$ gehört, so gehört auch jedes kleinere Element
von $\mathfrak R$ zu $A$. Ist $A$ ein Theil von $A'$, so wollen wir
$\alpha$ kleiner als $\alpha'$ nennen, und dadurch ist die Menge
$\mathfrak S$ zu einer geordneten geworden.

Sehen wir die durch die rationalen Brüche erzeugten Schnitte als
gleichwerthig mit diesen rationalen Brüchen selbst an und nennen sie
kurz \gesperrt{rationale Schnitte}, so enthält die Menge $\mathfrak S$
die Menge $\mathfrak R$, und $\mathfrak S$ ist also jedenfalls eine
\gesperrt{dichte Menge}. Die Menge $\mathfrak S$ ist aber auch
\gesperrt{stetig}; denn bezeichnen wir die Schnitte durch die grossen
deutschen Buchstaben $\mathfrak A,\mathfrak B,\ldots$ und betrachten
irgend einen Schnitt in der Mannigfaltigkeit der Schnitte,
$(\mathfrak A,\mathfrak B)$, so können wir ein Element in
$\mathfrak S$ bestimmen, $\alpha=(A,B)$, indem wir in $A$ jeden
rationalen Bruch aufnehmen, der einen der Schnitte von $\mathfrak A$
erzeugt, und alle anderen rationalen Brüche, die also die rationalen
Schnitte in $\mathfrak B$ erzeugen, nach $B$ werfen. Dieser Schnitt
$\alpha$ in $\mathfrak R$ erzeugt den Schnitt
$(\mathfrak A,\mathfrak B)$ in $\mathfrak S$. Dies wird nachgewiesen
sein, wenn gezeigt ist, dass jedes Element $\alpha'$ in $\mathfrak S$,
was kleiner ist als $\alpha$, zu $\mathfrak A$ gehört, und jedes
Element $\beta'$ in $\mathfrak S$, was grösser ist als $\alpha$, zu
$\mathfrak B$.

Sei also $\alpha'=(A',B')$ und $\alpha'<\alpha$, dann giebt es
rationale Brüche in $A$, die nicht in $A'$ enthalten sind; es giebt
also ein rationales $\mu$, so dass $\alpha'<\mu<\alpha$, und dieses
$\mu$ erzeugt einen Schnitt, der in $\mathfrak A$ enthalten ist,
gehört also selbst zu $\mathfrak A$; da $\alpha'<\mu$ ist, so gehört
auch $\alpha'$ zu $\mathfrak A$. Ganz ebenso zeigt man, dass jedes
$\beta'$, das grösser als $\alpha$ ist, zu $\mathfrak B$ gehört, und
damit ist die Stetigkeit von $\mathfrak S$ nachgewiesen.

Diese sehr abstracte Betrachtungsweise giebt uns die Sicherheit, dass
die Annahme einer stetigen Menge keinen Widerspruch enthält, dass
solche Mengen wenigstens im Reiche der Gedanken existiren. Die
Geometrie wie die Analysis, die immer gern an die geometrische
Anschauung anknüpft, hat lange stillschweigend die \gesperrt{Existenz}
stetiger Mengen, z.~B. bei den Punkten einer geraden Linie oder irgend
eines anderen zusammenhängenden Linienzuges, als eine Art von Axiom
angenommen. Auch der \gesperrt{Unterschied zwischen dichter und
stetiger Menge}, der der Unterscheidung commensurabler und
incommensurabler Strecken zu Grunde liegt, ist den Alten nicht
entgangen.\setcounter{footnote}{0}\footnote{Euklid, Elemente, Buch X.}
